% 增长的上限——逻辑斯蒂模型
% 与活动一对照：同一组人口数据、同一个“2050 年会怎样”的问题，这次给增长加上天花板。
% 活动十六用连续模型拟合人口，活动十七用离散差分模型讨论干预时机。

\section{活动十六：人口增长的天花板——逻辑斯蒂模型}

\subsection{情境}
回到活动一的城市：2000 年 800 万，2010 年 1000 万，2020 年 1250 万。
当时指数模型预测 2050 年约 2440 万。现在多了一条关键信息：该城的水资源、住房
与就业岗位大约只能支撑 1600 万人——环境容纳量 $K=1600$（万）。

\subsection{提出问题}
(1) 把“天花板”写进模型，重新拟合同一组数据；(2) 2050 年的预测变成多少？
(3) 活动一的指数预测从哪一年开始失真？

\subsection{选模型}
逻辑斯蒂模型（阻滞增长模型）认为：剩余空间越小，增长越慢。
\[
  \frac{dP}{dt}=rP\left(1-\frac{P}{K}\right),
\]
其解为
\[
  P(t)=\frac{K}{1+\left(\frac{K-P_0}{P_0}\right)e^{-rt}}.
\]
与活动一的指数模型对比：指数模型的相对增速恒为 $r$；逻辑斯蒂把它改成
$r\left(1-\frac{P}{K}\right)$。$P\ll K$ 时括号近似为 $1$，退化为指数增长；
$P\to K$ 时相对增速趋于 $0$，增长停止。

\subsection{拟合}
取 $K=1600$、$P_0=800$，则 $\dfrac{K-P_0}{P_0}=1$，模型简化为
$P(t)=\dfrac{1600}{1+e^{-rt}}$（万）。
用 2010 年的数据定 $r$：由 $P(10)=1000$ 得
$1+e^{-10r}=\dfrac{1600}{1000}=1.6$，即 $e^{-10r}=0.6$，所以
\[
  r=\frac{\ln(1/0.6)}{10}\approx0.0511.
\]
检验 2020 年：$P(20)=\dfrac{1600}{1+e^{-1.022}}\approx1176$（万），
比实际 1250 万低 74 万（约 $5.9\%$）。只用一个点定参，留出的检验点就暴露了偏差。

\subsection{预测：两条曲线分道扬镳}
\begin{center}
\footnotesize
\setlength{\tabcolsep}{5pt}
\renewcommand{\arraystretch}{1.3}
\begin{tabular}{lrrr}
\toprule
年份 & 实际 & 指数模型 & 逻辑斯蒂模型 \\
\midrule
2000 & 800 & 800 & 800 \\
2010 & 1000 & 1000 & 1000 \\
2020 & 1250 & 1250 & 1176 \\
2030 & --- & 1563 & 1316 \\
2040 & --- & 1953 & 1416 \\
2050 & --- & 2441 & 1485 \\
\bottomrule
\end{tabular}
\end{center}
（预测值由模型精确计算后四舍五入到万。）

数据段内两个模型都不差；差距在外推中迅速放大：2025 年相差约 146 万，
2030 年约 247 万，2050 年约 957 万。

指数模型何时开始失真？令 $800\times1.0226^{t}=1600$，得
$t=\dfrac{\ln2}{\ln1.0226}\approx31.1$，即 2031 年前后，指数预测首次超过天花板 $K$。
此后每外推一年，误差就多一分——2050 年的 2441 万比上限还多出 841 万，显然不可能。
逻辑斯蒂给出 1485 万，仍在上限之内，并继续缓慢逼近 1600 万。

机制：有效相对增速为 $r\left(1-\dfrac{P}{K}\right)$。$P=800$ 万时约 $0.026/\text{年}$，
$P=1000$ 万时约 $0.019/\text{年}$，$P=1250$ 万时只剩 $0.011/\text{年}$。
人口越多，剩余空间越小，增长越慢；指数模型把这一项永久设为 $1$，外推必然偏高。

\subsection{检验与反思}
逻辑斯蒂在 2020 年低估 74 万，说明真实增长比 $K=1600$、$r=0.0511$ 允许的更快一些
（也许实际容纳量更高，或存在人口迁入）。若对三个数据点做最小二乘，得 $r\approx0.0599$，
2020 年预测 1229 万、2050 年预测 1524 万——参数取值略有出入，结论不变。
指数模型在数据段内误差为零，恰是陷阱：它能完美复现历史，却会在物理上不可能的地方继续外推。
\key{指数模型描述“无限制复制”，逻辑斯蒂模型描述“有限空间里的增长”；外推越远，两者差距越大。}

\section{活动十七：传染病的暴发与干预时机——逻辑斯蒂差分模型}

\subsection{情境}
某校有 2000 名易感学生，开学初 5 人感染流感。学校卫生室力量有限，只能争取在某一周转入
强化措施（减少聚集、部分线上课），使每周的有效传播率下降。

\subsection{提出问题}
不干预，疫情多久见顶、最终多大？同样把传播率减半，在哪一周启动效果最好？

\subsection{选模型}
按周递推的离散逻辑斯蒂模型：
\[
  I_{n+1}=I_n+rI_n\left(1-\frac{I_n}{K}\right),\qquad K=2000,\quad I_0=5,\quad r=0.5,
\]
其中 $I_n$ 为第 $n$ 周的累计感染人数；每周结果四舍五入取整后再进入下一周。
每周新增 $=rI_n\left(1-\dfrac{I_n}{K}\right)$：感染基数 $I_n$ 越大、剩余易感空间
$\left(1-\dfrac{I_n}{K}\right)$ 越小，新增就越少。
初期 $I_n\ll K$ 时 $I_{n+1}\approx(1+r)I_n=1.5I_n$，即
$I_n\approx5\times1.5^{\,n}\approx5e^{0.406n}$，约每 $\dfrac{\ln2}{\ln1.5}\approx1.7$ 周翻一倍——
正是活动一指数模型的形式。

\subsection{计算：先看不干预}
\begin{center}
\footnotesize
\setlength{\tabcolsep}{4pt}
\renewcommand{\arraystretch}{1.25}
\begin{tabular}{rccr}
\toprule
周次 $n$ & 累计感染 $I_n$ & 单周新增 & 纯指数近似 $5\times1.5^{n}$ \\
\midrule
0 & 5 & --- & 5 \\
1 & 7 & 2 & 8 \\
2 & 10 & 3 & 11 \\
3 & 15 & 5 & 17 \\
4 & 22 & 7 & 25 \\
5 & 33 & 11 & 38 \\
6 & 49 & 16 & 57 \\
7 & 73 & 24 & 85 \\
8 & 108 & 35 & 128 \\
9 & 159 & 51 & 192 \\
10 & 232 & 73 & 288 \\
11 & 335 & 103 & 432 \\
12 & 474 & 139 & 649 \\
13 & 655 & 181 & 973 \\
14 & 875 & 220 & 1460 \\
15 & 1121 & 246 & 2189 \\
16 & 1367 & 246 & 3284 \\
\bottomrule
\end{tabular}
\end{center}

第 15–16 周新增最多（约 246 人），与理论峰值 $\dfrac{rK}{4}=250$ 接近；
第 20 周 1924 人，最终趋近 2000（全部易感者）。
指数近似前 3 周几乎重合，第 10 周已高估到 288（实际 232，高约 $24\%$）：
一旦 $I_n$ 不再远小于 $K$，指数近似立刻失效——与活动十六看到的失真同源。

\subsection{干预实验：什么时候动手}
设从第 $m$ 周起把 $r$ 降为 $0.2$ 并保持，重跑同一模型。
\begin{center}
\footnotesize
\setlength{\tabcolsep}{5pt}
\renewcommand{\arraystretch}{1.3}
\begin{tabular}{rccc}
\toprule
启动周 $m$ & 峰值单周新增 & 峰值出现周 & 第 16 周累计 \\
\midrule
0 & 100 & 33 & 86 \\
4 & 100 & 29 & 178 \\
8 & 100 & 24 & 402 \\
10 & 100 & 21 & 578 \\
11 & 103 & 11 & 685 \\
12 & 139 & 12 & 803 \\
13 & 181 & 13 & 933 \\
14 & 220 & 14 & 1073 \\
15 & 246 & 15 & 1220 \\
不干预 & 246 & 15 & 1367 \\
\bottomrule
\end{tabular}
\end{center}

规律很清晰：
\begin{itemize}
  \item 第 10 周之前启动，峰值都能压到约 100，正是新参数下的理论峰值
        $\dfrac{0.2\times2000}{4}=100$，峰值还被推迟到第 21–33 周；
  \item 第 12 周才启动，峰值回升到 139；第 13–14 周为 181、220；第 15 周之后基本无效。
\end{itemize}
决策准则：\textbf{赶在累计感染还不到 $K$ 的约十分之一时启动}（本例 $I_{10}=232$，
仍在近似指数段），等价于\textbf{赶在新增见顶前约 5 周}。道理在于从启动到新增被压住
要好几个星期；等新增自己放缓时再动手，大部分感染已经发生。

\subsection{检验与反思}
注意：干预只改变 $r$，并没有改变 $K$——各种情形最终都趋近 2000。
降低 $r$ 买到了时间、压低了峰值，却没有减少总感染人数。
要真正减少总感染人数，必须减小 $K$ 本身：接种疫苗把易感者移出、隔离传染源。
另外，模型假设措施长期维持；若几周后放松，$r$ 恢复，曲线会重新加速。
\key{逻辑斯蒂模型里，$K$ 决定最终规模，$r$ 决定到达的速度：压峰靠 $r$，减总量靠 $K$。}
